\section{验证缺口与“思想未经严格推演”}

第 4 章记录了“没有 Lean、没有评审”这一事实。本章往前多走一步：把“验证缺口”拆成可以分别追问的三层，并给出一条可执行的核验路线。需要先把分寸立住——本文说“缺口”，说的是证据缺口与理解缺口，不是判定手稿有错。这两件事在数学共同体的日常语用里经常被混为一谈，而它们对研究决策的含义完全不同。

\subsection{三层缺口}

\textbf{第一层：依赖链未验证。}四维篇把自己的启动成本压在了三维姊妹篇的加权板条引理上，后者又立在 Guth--Wang--Zahl 的并集估计上（第 4.3 节）。链条本身没有可疑之处，但意味着验证必须从链条根部做起：先确认三维引理段的推导，再回到四维篇的归约。反过来，如果只核验四维篇第 10 章的收口段，等于在一根未检的承重梁上做局部检修。这类结构性依赖在人类数学里很常见（一篇文章引用另一篇的引理），它之所以在此成为风险，是因为两篇稿件的验证状态同为零，而验证它们的预期难度又都远高于常规文章。

\textbf{第二层：概念未蒸馏。}地平线、两色匹配、水平模型这些自造词汇，目前只在一份手稿里被使用过一次。数学史上，一个概念要“立住”，需要两样东西：别处的独立使用，以及教学意义上的可复述性。以粘性（sticky）为例：Katz--\L aba--Tao 在 2000 年把它作为极端形态的标签提出\upcite{klt2000}，之后二十余年里，它先是变成 Wang--Zahl 的独立定理对象\upcite{wangz2022}，又出现在 Furstenberg 问题、Nikodym 转移与弯曲挂谷的文献中\upcite{wangzakharov2026,gaoliuxi2025,raichoudhuri2026}——概念的跨问题复用，是它被共同体真正消化的实证。地平线会不会走上同一条路，现在无从判断；它甚至连“由人类数学家在讨论班上重新讲一遍”这一步都还没有发生。

\textbf{第三层：推演未受检。}175 页的论证是一个自足的闭环，闭环的正确性只能靠逐段复演来确认。稿件中至少有两类步骤属于“机制最密、最容易藏下隐性假设”的位置：一是比较原理里“导数失配 → 多项式端点映射体积超界”的推导（依赖 Lem\-ma 3.3 型的导数、采样界），二是收口段把局部改进“回填”到原问题的转换（依赖基网格全变差控制）。这两处的共同特点是：如果某条技术性假设过窄，错误不会以“公式矛盾”的方式暴露，而会以“某条不等式在某个子情形下差一个幂次”的方式潜伏——这恰恰是人类审读（以及形式化）最需要耗时的地方。

\begin{keybox}[三层缺口的共同点]
三层缺口指向的都是“证据还没产生”，不是“证据已经反转”。当前对 074 族的正确判断是“未知”，而不是“可疑”——但“未知”在验证成本上等价于一项尚未开工的大型工程，这一点常被讨论忽略。
\end{keybox}

\subsection{“证明”与“理解”是两种成就}

数学共同体对一个结果的接受，实际包含三件事：可核验（有专家能独立复演）、可理解（核心思想能被讲授与改造）、可继承（概念能被其他问题复用）。手稿目前只交付了第四件东西：一份格式完整的文本。

AHM 声明的落点正在“理解”一侧：其结尾呼吁“回到以人类理解为中心的科学愿景”\upcite{ahm2026}。这个立场与“结果是否正确”正交——即使手稿完全正确，理解缺口依然存在。反过来说，也不应把“理解缺口”当作“结果无效”的证据：Newton 时代的许多证明，其可理解性同样是后世几代人才补上的。本文的立场是把两笔账分开记：正确性账（技术层，第 9.4 节的路线图服务于它）与理解账（概念层，需要教学、综述与再表述工作，第 10 章第 3 轨服务于它）。

\subsection{Lean 形式化能覆盖什么、覆盖不了什么}

形式化在此案中的角色需要精确表述，既不能高估也不能低估。它能做的：把一段论证翻译成机器可检查的项，消除“隐性假设”“手滑的常数”“错误的引用编号”这类机械错误，并且把核验的工作量从“全文重读”压缩为“读形式化脚本”。它做不到的三件事：其一，形式化验证的对象是翻译后的命题，翻译本身的保真度不受保护——近期有专门研究以 Navier--Stokes 的案例说明“自动形式化通过检验”与“自然语言证明正确”之间存在系统性缝隙\upcite{bastounis2026}；其二，形式化不判断结果的新意与概念价值；其三，形式化不产生“理解”。

因此，074 族“未形式化”的含义要这样读：不是“有多一层证书”与“没有证书”的差别，而是“机械核验尚未开工，且该领域已经证明这类开工是可行的”——三维挂谷证明（王虹--Zahl）的 Lean 4 形式化已由中科院 MMAT 与 Project Numina 等合作完成\upcite{numina2026}。同一题材、相近量级，人类侧做到了，074 族没有做。为什么没有做，公开材料没有说明；从工作量与稿件整理状态推测（\cfL），最可能的原因不是“不可能”，而是“没有排进优先级”。

\subsection{可复现核验路线图}

\begin{table}[t]
\centering
\caption{四维稿件的模块化核验路线（编号 M1--M8）}
\label{tab:modules}
\footnotesize
\begin{tabularx}{\textwidth}{@{}L{0.9cm}L{3.4cm}L{1.2cm}Y@{}}
\toprule
\rowcolor{lightbg}\textbf{模块} & \textbf{内容} & \textbf{体量} & \textbf{核验要点与风险} \\
\midrule
M1 & 归约（Theorem 2.17 与第 2 章机件） & 约 20 页 & 决定论证是否启动；核对无测度性假设下每步选择；容量类 $\mathcal{C}(d)$ 两条条件的导出 \\
M2 & 比较原理（第 3 章） & 约 22 页 & 导数/采样界（Lemma 3.3 型）是全部列表匹配的前提；体积超界论证的计数 \\
M3 & 临界路径与改进转换（第 4 章） & 约 24 页 & Proposition 4.11 与 Lemma 4.12 的“回填”；指数约定的全局一致性 \\
M4 & 各向同性（第 5--6 章） & 约 35 页 & 次数恢复（Proposition 6.8）的 Remez/Vandermonde 常数控制；圆锥障碍的完备性 \\
M5 & 保边平面刚性（第 7 章） & 约 14 页 & Proposition 7.6 的“保留原始边”与乘法卷积估计的接口 \\
M6 & 有限窄度与水平模型（第 8--9 章） & 约 37 页 & 临界速率 $2k=1$ 的三腿切换；$2k<1$ 自举的终止性 \\
M7 & 无界窄度与收口（第 10 章） & 约 15 页 & 密度指数漂移在比值中的抵消；最终矛盾的指数差为正 \\
M8 & 几何推论（第 11 章） & 约 7 页 & 转移定理假设（Keleti--M\'ath\'e、Gao--Liu--Xi、Nadjimzadah）逐条对照 \\
\bottomrule
\end{tabularx}
\end{table}

建议顺序与理由：先 M1，再“链条根”（三维篇中加权板条引理所在的一节，结合 Guth--Wang--Zahl 的并集估计原文），然后 M2、M3。这四块合计约 70 页，却是整条链上唯一“不可绕开”的部分——其余模块可以按人力并行。形式化的优先级与之类似：M1 与三维引理段的结构最规整、自足性最强，局部形式化的性价比最高；三区间（M4--M7）体量大且依赖分析机件，放在靠后。

\begin{dirbox}[独立小课题（可作为核验的副产品）]
（i）把三维篇 Lemma 2.3 的“加权凸聚类 → 加权板条界”推导整理成人可读版本，作为核验附记；（ii）对 M1 写一份逐步假设审计（audit），逐行确认没有引入可测性、紧性假设；（iii）对三腿切换与锥障碍构造有限离散模型（如 $\mathbb{F}_p$ 类比或小规模数值实验）做压力测试；（iv）对 Proposition 6.8 尝试交互式证明助手内的局部形式化；（v）对照转移定理假设，核对第 11.4--11.5 节推论的条件表。每条产出独立、可验收，见附录 B 的验收标准建议。
\end{dirbox}

\subsection{为什么“正确”之外还要追问“推演”}

写到这里可以给出本章的收束判断：对 074 族的讨论，最好按“三张账”分开记——正确性账（未开工，但有明确路线）、理解账（未开工，需要概念蒸馏与再表述）、影响账（争议中，取决于共同体如何消化）。把三张账混成一张，就会得到两种常见但都站不住的结论：或者因为“AI 生成”而预判不可信，或者因为“结论漂亮”而预支信任。两种结论的共同错误，是让立场代替证据。

\begin{keybox}[可记忆主句]
手稿最稀缺的东西不是篇幅，而是“第二次使用”：一旦地平线、图卡与比较原理在别人的论文里被独立使用一次，这份工作才算真正进入数学的毛细血管。
\end{keybox}